\documentclass[11pt]{article}
\usepackage[review]{acl}
\usepackage{times}
\usepackage{latexsym}
\usepackage[T1]{fontenc}
\usepackage[utf8]{inputenc}
\usepackage{microtype}
\usepackage{booktabs}
\usepackage{graphicx}
\usepackage{amsmath}

\newcommand{\best}[1]{\textbf{#1}}

\title{Prompt-Injection Detectors Do Not Transfer, \\ but a LoRA-Tuned LLM Generalizes:
A Held-Out Cross-Dataset Study}

\author{Anonymous ACL submission}

\begin{document}
\maketitle

\begin{abstract}
Leading prompt-injection detectors report near-perfect recall---ProtectAI's classifier 99.7\%,
Meta Prompt-Guard 99.5\% TPR---yet those numbers are almost always measured \emph{in-distribution}
on the detector's own dataset. We ask whether they survive a
change of dataset. Evaluating three detectors---a regex pre-filter, a fine-tuned DeBERTa classifier
(ProtectAI), and an 8B LLM tuned with a single LoRA adapter---on a strict held-out protocol (train
on one public injection set, test on two others never seen in training), we find that reported
recall does not transfer: ProtectAI, which reports 99.7\% in-distribution, drops to 0.40 recall on a
held-out set. A regex catalogue is worse still (0.20). The LoRA-tuned LLM, trained only on the same
single set, \emph{generalizes}: it reaches 0.97 ROC-AUC on \emph{both} held-out benchmarks
(0.59 / 0.60 recall at 1\% FPR), exceeding the fine-tuned classifier on deepset (0.88 / 0.40) and
matching its AUC on jailbreak, where reported detectors already do well. The generalization comes
with a cost we report honestly: the LLM detector is \emph{more} vulnerable to trivial character
perturbations (14\% zero-width evasion vs.\ 0\% for the classifier). Our findings argue that
detection recall is a dataset-fragile quantity and that no single detector---however strong
in-distribution---should be treated as a security guarantee; it is one evadable tier of a
defense-in-depth stack.
\end{abstract}

\section{Introduction}
Indirect prompt injection---malicious instructions smuggled into the data an LLM or LLM agent
consumes---is the top-ranked risk in the 2025 OWASP Top~10 for LLM
Applications~\citep{owasp2025}. A common first line of defense is a \emph{detector}: a classifier
that flags injection-bearing inputs before the model acts on them. Commercial and open detectors
report striking numbers---ProtectAI's DeBERTa-v3 injection classifier reports 99.7\%
recall~\citep{protectai}, and Meta's Prompt-Guard reports 99.5\% TPR at AUC
$>$0.99~\citep{promptguard}.

These headline numbers are, almost without exception, measured \emph{in-distribution}: the test set
is drawn from the same distribution (often the same corpus) as the training set. Deployment is
out-of-distribution by definition---the attacks a detector meets in production are not the ones in
its training set. We therefore ask a simple question: \emph{do reported detector numbers survive a
change of dataset?}

We evaluate three detectors spanning the practical design space---a hand-written regex pre-filter, a
fine-tuned encoder classifier (ProtectAI DeBERTa-v3), and an 8B decoder LLM adapted with a single
LoRA~\citep{hu2022lora} adapter---under a strict held-out protocol: all detectors are scored on the
\emph{identical} rows of two public injection benchmarks that the trained detectors never saw in
training. Our contributions:
\begin{enumerate}\setlength\itemsep{2pt}
\item A controlled cross-dataset evaluation showing that reported in-distribution recall
does not transfer: the fine-tuned classifier's recall falls from a reported 99.7\% to
\best{0.40} on a held-out set (\S\ref{sec:results}).
\item A LoRA-tuned LLM detector that, trained on a single injection set, \emph{generalizes
uniformly} (0.97 AUC on both held-out sets) and exceeds the fine-tuned classifier on the set where
the classifier fails to transfer.
\item An honest characterization of the accompanying trade-off: the LLM detector generalizes
better but is \emph{more} vulnerable to trivial character-level evasion.
\end{enumerate}
The practical message is not ``use an LLM detector''; it is that detection recall is fragile
across datasets and evadable under perturbation, so a detector should be one defense-in-depth tier,
never the guarantee.

\section{Related Work}
\paragraph{Prompt-injection detection.}
Detection-based defenses classify the input text. Open detectors include ProtectAI's
DeBERTa-v3-base classifier~\citep{protectai}, Meta Prompt-Guard~\citep{promptguard}, and
LLM-as-judge preprocessors~\citep{promptarmor}. Benchmarks such as
PINT~\citep{pint} and the over-defense benchmark NotInject~\citep{injecguard} rank detectors, but
each detector's headline recall is reported on its own evaluation split. Recent work shows detectors
collapse under adversarial character perturbation~\citep{guardrails}, and that firewall-style
detection is insufficient against adaptive attacks~\citep{firewalls}.

\paragraph{Design-level defenses.}
Because detection is evadable, a parallel line moves the guarantee to the \emph{action}:
CaMeL~\citep{camel} enforces value-level information-flow policy, Progent~\citep{progent} attaches
programmable privilege control to tool calls, and FIDES~\citep{fides} applies deterministic
information-flow control. These treat detection as at most a non-authoritative pre-filter---the
stance our results support empirically.

\paragraph{Parameter-efficient tuning.}
LoRA~\citep{hu2022lora} adapts a frozen base model with low-rank updates, enabling cheap
task-specialized adapters. We use it to turn a general 8B LLM into an injection detector and study
whether the base model's breadth yields better cross-dataset generalization than a small fine-tuned
encoder.

\section{Experimental Setup}\label{sec:setup}
\paragraph{Datasets.} We use three public prompt-injection classification sets:
\texttt{safe-guard-prompt-injection} (10{,}296 rows), \texttt{deepset/prompt-injections} (662), and
\texttt{jackhhao/jailbreak-classification} (1{,}306). Each row is labeled injection or benign.

\paragraph{Held-out protocol.} The trained LLM detector is fine-tuned \emph{only} on
\texttt{safe-guard} (a 95/5 train/val split). \texttt{deepset} and \texttt{jailbreak} are
\emph{held out entirely}---never seen in training---and used only for evaluation. This isolates
out-of-distribution generalization: high held-out recall reflects transfer, not memorization. All
detectors are scored on the identical held-out rows.

\paragraph{Detectors.}
\emph{(i)~Regex pre-filter}: a catalogue of $\sim$100 hand-written patterns for known
injection phrasings, scored by a risk value.
\emph{(ii)~Fine-tuned classifier}: \texttt{protectai/deberta-v3-base-prompt-injection-v2}
(184M params), a sequence classifier we run locally on the same rows.
\emph{(iii)~LoRA-tuned LLM}: Qwen3-8B~\citep{qwen3} with a rank-16 LoRA adapter (attention $+$ MLP
projections), trained for one epoch with a completion-only objective to emit \texttt{INJECTION}
or \texttt{SAFE}; at inference we score the log-probability margin between the two labels. Training
used bf16 on a single 24\,GB GPU; the adapter is 175\,MB.

\paragraph{Metrics.} We report ROC-AUC and recall at a fixed low false-positive operating point,
matching how a production guardrail is judged. To probe robustness we apply three
character-perturbation transforms to caught positives---space-out (spaces between characters),
zero-width insertion, and leetspeak---and report the \emph{evasion rate}: the fraction of previously
caught injections that flip to missed after the transform.

\section{Results}\label{sec:results}
\paragraph{Reported recall does not transfer.}
Table~\ref{tab:main} shows the held-out results. The fine-tuned classifier, which reports
99.7\% recall on its own held-out data~\citep{protectai}, achieves only \best{0.40} recall on
\texttt{deepset}: a change of dataset erases most of its measured recall. The regex catalogue is
weaker still (0.20)---its fixed patterns miss the persona-injection, role-play, and task-pivot
phrasings that dominate \texttt{deepset} and carry no lexical trigger.

\paragraph{The LoRA-tuned LLM generalizes.}
Trained only on \texttt{safe-guard}, the LoRA detector reaches \best{0.97} ROC-AUC on
\emph{both} held-out sets. On \texttt{deepset} it exceeds the fine-tuned classifier on both AUC
(0.97 vs.\ 0.88) and recall (0.59 vs.\ 0.40) on data neither model saw in training---the base
model's breadth transfers across injection distributions where a small task-specific encoder does
not. On \texttt{jailbreak} the two are closer: the LoRA matches ProtectAI's AUC (0.97 vs.\ 0.98) but
ProtectAI retains a higher operating-point recall (0.89 vs.\ 0.60), a set where the encoder
already generalizes well. The consistent reading is that the LoRA detector is uniformly strong
across distributions (0.97 AUC on both), whereas the encoder is strong on one held-out set (0.98)
and weak on the other (0.88).

\paragraph{Generalization trades off against evasion robustness.}
The gain is not free (Table~\ref{tab:evasion}). Under character perturbation the LLM detector is
\emph{more} evadable than the fine-tuned classifier: 14\% of caught \texttt{deepset} injections
evade it under zero-width insertion (vs.\ 0\% for ProtectAI), and 5\% under space-out; on
\texttt{jailbreak} the zero-width evasion rises to 37\%. The regex
pre-filter, whose normalization strips these transforms, is near-0\% evadable---but it caught little
to begin with. Higher recall and higher perturbation-robustness are, on this axis, in tension.

\begin{table}[t]
\centering\small
\setlength\tabcolsep{5pt}
\caption{Held-out detection (never trained on these sets). ROC-AUC / recall at 1\% FPR.
Best per column in \best{bold}. ProtectAI's \emph{reported} in-distribution recall is 99.7\%.}
\label{tab:main}
\begin{tabular}{lcc}
\toprule
Detector & deepset (662) & jailbreak (1{,}306) \\
\midrule
Regex pre-filter & 0.60 / 0.20 & 0.75 / 0.35 \\
ProtectAI DeBERTa-v3 & 0.88 / 0.40 & \best{0.98} / \best{0.89} \\
LoRA-tuned Qwen3-8B & \best{0.97} / \best{0.59} & \best{0.97} / 0.60 \\
\bottomrule
\end{tabular}
\end{table}

\begin{table}[t]
\centering\small
\setlength\tabcolsep{4pt}
\caption{Character-perturbation evasion rate on held-out \texttt{deepset} (lower is better):
fraction of caught injections that flip to missed after the transform.}
\label{tab:evasion}
\begin{tabular}{lccc}
\toprule
Detector & space-out & zero-width & leetspeak \\
\midrule
Regex pre-filter & \best{0.00} & \best{0.00} & \best{0.00} \\
ProtectAI DeBERTa-v3 & \best{0.00} & \best{0.00} & 0.08 \\
LoRA-tuned Qwen3-8B & 0.05 & 0.14 & 0.03 \\
\bottomrule
\end{tabular}
\end{table}

\section{Discussion}
Two findings have direct deployment consequences. First, \emph{detector recall is a
dataset-fragile quantity}: a number measured on one corpus is not a promise about another, so
procurement and red-teaming should evaluate detectors on held-out distributions, not vendor splits.
Second, \emph{generalization and evasion-robustness pull in opposite directions}: the detector that
transfers best is the one most easily perturbed, and the detector most robust to perturbation
(the regex) transfers worst. There is no free lunch that a single detector resolves.

Together these argue against resting a security guarantee on any detector. The complementary line of
design-level defenses~\citep{camel,progent,fides} that gate the \emph{action} by input provenance,
rather than classifying the \emph{text}, is attractive precisely because it does not inherit the
detector's dataset-fragility or perturbation-fragility. A trained detector like ours is best used as
a strong-but-non-authoritative pre-filter feeding such a boundary.

\section{Conclusion}
Reported prompt-injection detector numbers are measured in-distribution and do not transfer: a
classifier reporting 99.7\% recall falls to 0.41 on a held-out set. A LoRA-tuned 8B LLM,
trained on a single set, generalizes better across datasets---but is more evadable under trivial
perturbation. Detection is a fragile, evadable tier; the guarantee belongs at the action.

\section*{Limitations}
Our study covers three English text-classification injection sets and two trained detectors; it does
not cover multilingual, multimodal, or agentic (tool-call) injection, nor adaptive attacks that
optimize against a specific detector. The LoRA detector is trained for a single epoch on one set;
more data or epochs could shift both its generalization and its robustness. The held-out
generalization result is demonstrated on two sets---\texttt{jailbreak} results in
Table~\ref{tab:main} should be read alongside \texttt{deepset}, and broader held-out coverage is
future work. Our evasion probe uses three fixed character transforms and is a lower bound on
adversarial vulnerability, not a worst case. Finally, detector scores are not directly comparable to
end-to-end agent attack-success rate; a caught input is prevention only if the downstream system
acts on the detector's verdict.

\section*{Ethics Statement}
This work studies \emph{defenses}. All datasets are public prompt-injection classification
benchmarks; we introduce no new attack technique and release only detector evaluation code. Improved
injection detection is dual-use in the narrow sense that any robustness evaluation informs attackers,
but the net effect of clarifying that detector numbers do not transfer is to discourage
over-reliance on a single evadable defense.

\bibliography{references}
\bibliographystyle{acl_natbib}

\end{document}
